\chapter{โมเมนตัมและการดลในระบบกายภาพ}
\label{ch:ch05_momentum_impulse}
\index{โมเมนตัมและการดล}

\begin{conceptbox}[title=วัตถุประสงค์การเรียนรู้ประจำบท]
เมื่อสิ้นสุดการศึกษาบทเรียนนี้ นักศึกษาสามารถ
\begin{enumerate}
    \item เมื่อศึกษาบทนี้จบแล้ว ผู้เรียนสามารถ
    \item นิยามโมเมนตัมเชิงเส้นและการดล พร้อมทั้งอธิบายความสัมพันธ์ระหว่างปริมาณทั้งสองผ่านทฤษฎีบทการดล-โมเมนตัม
    \item คำนวณโมเมนตัมของวัตถุเดี่ยวและของระบบวัตถุหลายชิ้น รวมทั้งคำนวณการดลจากแรงและช่วงเวลาที่แรงกระทำ
    \item อธิบายกฎการอนุรักษ์โมเมนตัมและระบุเงื่อนไขที่ทำให้โมเมนตัมของระบบคงตัว
    \item ใช้กฎการอนุรักษ์โมเมนตัมแก้โจทย์ปัญหาการชนได้อย่างเป็นระบบ
\end{enumerate}
\end{conceptbox}

\section{ทฤษฎีและหลักการพื้นฐาน}

\begin{bioappbox}
\textbf{การถ่ายโอนโมเมนตัมในการแพร่ระดับโมเลกุลและความดันออสโมซิส}

ในระดับโมเลกุล การชนกันของโมเลกุลน้ำและสารละลายกับผนังเซลล์ทำให้เกิดการถ่ายโอนโมเมนตัม $\Delta p = F\Delta t$ ซึ่งเป็นที่มาของความดันทางจลน์ศาสตร์ (Kinetic Pressure) และแรงดันออสโมซิส (Osmotic Pressure) ตามสมการของฟานต์ฮอฟฟ์ $\Pi = iCRT$ ความเข้าใจเรื่องโมเมนตัมยังช่วยอธิบายการกระจายตัวของอนุภาคละอองลอย (Bioaerosols) เมื่อเกิดการจามหรือการไอในสภาพแวดล้อมปิด
\end{bioappbox}

\begin{labbox}
1. ตรวจสอบตำแหน่งศูนย์ (Zero Error) ของเครื่องมือวัดทุกชนิดก่อนเริ่มบันทึกข้อมูล
2. ทำการวัดซ้ำอย่างน้อย 3 ถึง 5 ครั้งในแต่ละจุดทดลองเพื่อคำนวณหาค่าเฉลี่ยและค่าเบี่ยงเบนมาตรฐาน
3. บันทึกผลการวัดพร้อมระบุหน่วยเอสไอ (SI Units) และค่าความคลาดเคลื่อนของการวัดเสมอ
\end{labbox}

\section{ตัวอย่างโจทย์และการวิเคราะห์เชิงฟิสิกส์}

\begin{examplebox}
1) การชนแบบยืดหยุ่น (elastic collision) เป็นการชนที่พลังงานจลน์รวมของระบบก่อนและหลังการชนมีค่าเท่ากัน กล่าวคือ ทั้งโมเมนตัมและพลังงานจลน์ต่างก็คงตัว การชนประเภทนี้ในธรรมชาติพบได้ชัดเจนที่ระดับอนุภาคหรืออะตอม (เช่น การชนกันของโมเลกุลก๊าซ) ส่วนในระดับมหภาค การชนของลูกบิลเลียดหรือลูกแก้วที่มีความแข็งมากถือเป็นตัวอย่างที่ใกล้เคียงกับการชนแบบยืดหยุ่นในอุดมคติ เมื่อใช้กฎการอนุรักษ์โมเมนตัมร่วมกับกฎการอนุรักษ์พลังงานจลน์ จะสามารถแก้สมการหาความเร็วหลังชนของวัตถุทั้งสองได้ดังนี้

v₁' = [(m₁-m₂)/(m₁+m₂)]v₁ + [2m₂/(m₁+m₂)]v₂

v₂' = [2m₁/(m₁+m₂)]v₁ + [(m₂-m₁)/(m₁+m₂)]v₂

2) การชนแบบไม่ยืดหยุ่น (inelastic collision) เป็นการชนที่พลังงานจลน์รวมของระบบหลังการชนมีค่าน้อยกว่าก่อนการชน แม้โมเมนตัมยังคงอนุรักษ์อยู่เสมอ พลังงานจลน์ส่วนที่หายไปจะถูกเปลี่ยนรูปไปเป็นพลังงานรูปแบบอื่น เช่น ความร้อน เสียง หรือพลังงานที่ใช้ในการเปลี่ยนรูปร่างของวัตถุ (การยุบตัว การเสียดสีภายในวัสดุ) การชนในชีวิตประจำวันส่วนใหญ่ รวมถึงอุบัติเหตุทางรถยนต์ จัดอยู่ในประเภทนี้
\end{examplebox}

\section{แบบฝึกหัดทบทวนและคำถามท้ายบท}

\begin{enumerate}
    \item จงอธิบายความแตกต่างระหว่างโมเมนตัมและการดล พร้อมระบุหน่วยของแต่ละปริมาณ
    \item เพราะเหตุใดโมเมนตัมจึงถูกจัดเป็นปริมาณเวกเตอร์ จงยกตัวอย่างประกอบ
    \item รถยนต์มวล 1,200 kg เคลื่อนที่ด้วยความเร็ว 20 m/s จงหาขนาดของโมเมนตัมของรถยนต์
    \item ลูกเทนนิสมวล 0.06 kg ถูกตีจนความเร็วเปลี่ยนจาก 10 m/s เป็น -20 m/s (ทิศตรงข้าม) ถ้าไม้เทนนิสสัมผัสลูกบอลเป็นเวลา 0.01 s จงหาแรงเฉลี่ยที่ไม้เทนนิสกระทำต่อลูกบอล
    \item จงอธิบายว่าเหตุใดการยืดเวลาการปะทะจึงช่วยลดแรงกระแทกที่วัตถุได้รับ โดยอ้างอิงทฤษฎีบทการดล-โมเมนตัม
    \item จงระบุเงื่อนไขที่ทำให้กฎการอนุรักษ์โมเมนตัมสามารถนำมาใช้ได้กับระบบหนึ่ง ๆ
\end{enumerate}

% สิ้นสุดบทเรียน
